% !TEX root = task-6-9-2.tex
% Editable task bank. Append new BankTask blocks at the end of the document.
\documentclass[12pt,a4paper]{article}
\usepackage{physics-textbook}
\usepackage{graphicx,needspace}
\usetikzlibrary{decorations.pathmorphing}
\newif\ifShowSolutions
\ifdefined\TutorEdition
  \ShowSolutionstrue
\else
  \ShowSolutionsfalse % Change to \ShowSolutionstrue to display solutions.
\fi
\newcounter{banktask}
\newsavebox{\BankTaskBox}
\newcommand{\BankTaskContent}[3]{%
  \noindent\textbf{\thebanktask.}\quad #2\par
  \ifShowSolutions
    \Needspace{50mm}% Keep the solution heading with its opening text or diagram.
    #3\par
    \ifdefined\TutorEdition
      \ifcsname TutorSource#1\endcsname
        \nopagebreak\smallskip{\small\csname TutorSource#1\endcsname\par}%
      \fi
    \fi
  \fi
}
\newcommand{\BankTask}[3]{%
  \refstepcounter{banktask}\par\medskip
  \sbox{\BankTaskBox}{\begin{minipage}{\linewidth}%
    \BankTaskContent{#1}{#2}{#3}\end{minipage}}%
  \ifdim\dimexpr\ht\BankTaskBox+\dp\BankTaskBox\relax>\dimexpr\textheight-20mm\relax
    \Needspace{35mm}\BankTaskContent{#1}{#2}{#3}%
  \else
    \Needspace{\dimexpr\ht\BankTaskBox+\dp\BankTaskBox\relax}%
    \noindent\usebox{\BankTaskBox}\par
  \fi
  \medskip
}
\begin{document}
\BookSubsection{6.9.2}{Радиоактивные превращения}


% Source sdamgia-2225: https://phys-ege.sdamgia.ru/problem?id=2225
\BankTask{sdamgia-2225}{
В результате нескольких $\alpha $- и $\beta $-распадов ядро урана ${}^{238}_{92}$U превращается в ядро свинца ${}^{206}_{82}$Pb. Определите количество $\alpha $-распадов и количество $\beta $-распадов в этой реакции.

\begin{center}\small\renewcommand{\arraystretch}{1.25}\begin{tabular}{|p{0.450\linewidth}|p{0.450\linewidth}|}\hline
Количество $\alpha $-распадов & Количество $\beta $-распадов\\\hline
 & \\\hline
\end{tabular}\end{center}
}{
\smallskip\noindent\textbf{Решение}\par
Каждый $\alpha $-распад приводит к уменьшению массового числа ядра на $4$. Каждый $\beta $-распад не изменяет массовое. Следовательно, для того чтобы определить число $\alpha $-распадов, достаточно проследить изменение массового числа. При превращении ядра ${}^{238}_{92}\mathrm{U}$ в ядро ${}^{206}_{82}\mathrm{Pb}$ массовое число уменьшается на $32$. Таким образом, число $\alpha $-распадов равно $\frac{32}{4} =8.$ 

Определим количество $\beta $-распадов. После восьми $\alpha $-распадов заряд ядра станет равен $92 - 8 \cdot  2=76.$ Следовательно, чтобы в результате получилось ядро ${}^{206}_{82}\mathrm{Pb},$ необходимо, чтобы произошло еще $82 - 76=6$ $\beta $-распадов, которые и «подправят» зарядовое число.

Ответ: $86$.
}

% Source sdamgia-2319: https://phys-ege.sdamgia.ru/problem?id=2319
\BankTask{sdamgia-2319}{
Изотоп ${}^{226}_{88}\mathrm{Ra}$ превратился в изотоп ${}^{206}_{82}\mathrm{Pb}.$ При этом произошло X $\alpha $-распадов и Y $\beta $-распадов.

Чему равны X и Y?

\begin{center}\small\renewcommand{\arraystretch}{1.25}\begin{tabular}{|p{0.450\linewidth}|p{0.450\linewidth}|}\hline
X & Y\\\hline
 & \\\hline
\end{tabular}\end{center}
}{
\smallskip\noindent\textbf{Решение}\par
Каждый $\alpha $-распад приводит к уменьшению зарядового числа ядра на $2$ и уменьшению массового числа на $4$. Каждый $\beta $-распад не изменяет массовое число ядра и увеличивает зарядовое число ядра на $1$. Поскольку массовое число при превращении ${}^{226}_{88}\mathrm{Ra}$ в ${}^{206}_{82}\mathrm{Pb}$ изменилось на $226 - 206=20,$ заключаем, что произошло $\frac{20}{4} =5$ $\alpha $-распадов. Изменение зарядового числа на $88 - 82=6$ означает, что помимо $5$ $\alpha $-распадов произошло также $5 \cdot  2 - 6=4$ $\beta $-распада.

Ответ: $54$.
}

% Source sdamgia-7796: https://phys-ege.sdamgia.ru/problem?id=7796
\BankTask{sdamgia-7796}{
В результате некоторого числа $\alpha $-распадов и некоторого числа электронных $\beta $-распадов из ядра ${}^{A}_{Z}\mathrm{X}$ получается ядро $^{A - 8} _Z - 1 Y.$ Чему равно число $\beta $-распадов в этой ядерной реакции?
}{
\smallskip\noindent\textbf{Решение}\par
При ядерных реакциях сохраняются суммарные массовые и зарядовые числа частиц, вступивших в реакцию. Массовое число ядра уменьшилось на $8$, а зарядовое — на единицу. Следовательно, ядро испытало два $\alpha $-распада и три $\beta $-распада.

Ответ: $3$.
}

% Source sdamgia-8721: https://phys-ege.sdamgia.ru/problem?id=8721
\BankTask{sdamgia-8721}{
В результате серии радиоактивных распадов ядро урана ${}^{238}_{92}\mathrm{U}$ превращается в ядро свинца ${}^{206}_{82}\mathrm{Pb}.$ Какое количество $\alpha $- и $\beta $-распадов оно испытывает при этом?

\begin{center}\small\renewcommand{\arraystretch}{1.25}\begin{tabular}{|p{0.450\linewidth}|p{0.450\linewidth}|}\hline
Количество $\alpha $-распадов & Количество $\beta $-распадов\\\hline
 & \\\hline
\end{tabular}\end{center}
}{
\smallskip\noindent\textbf{Решение}\par
При альфа-распаде из ядра вылетает ядро гелия с массой $4$ и зарядом $2$, в результате чего заряд ядра уменьшается на два, а массовое число — на четыре единицы. При бета-распаде из нейтрон в ядре распадается на электрон и протон, в результате чего заряд ядра атома увеличивается на единицу, при этом масса ядра не изменяется.

При превращении ядра урана в ядро свинца масса уменьшается на $238 - 206 = 32$ а. е. м. Значит, происходит $32$ : $4 = 8$ альфа-распадов.

За эти $8$ альфа-распадов заряд ядра уменьшается на $2$ · $8 = 16$, а поскольку при превращении ядра урана в ядро свинца масса заряд уменьшается только на $92 - 82 = 10$ зарядов электрона, значит, происходит $16 - 10 = 6$ бета-распадов.

Ответ: $86$.
}

% Source sdamgia-10952: https://phys-ege.sdamgia.ru/problem?id=10952
\BankTask{sdamgia-10952}{
Ядро ${}^{20}_{12}\mathrm{Mg}$ испустило протон, а затем захватило электрон. Сколько протонов и сколько нейтронов входит в состав ядра, которое образовалось в результате этих реакций?

\begin{center}\small\renewcommand{\arraystretch}{1.25}\begin{tabular}{|p{0.450\linewidth}|p{0.450\linewidth}|}\hline
Число протонов & Число нейтронов\\\hline
 & \\\hline
\end{tabular}\end{center}

В ответе запишите число протонов и нейтронов слитно без пробела.
}{
\smallskip\noindent\textbf{Решение}\par
В результате такой реакции заряд ядра уменьшается на $2$, а массовое число уменьшается на $1$. У ${}^{20}_{12}\mathrm{Mg}$ заряд ядра равен $12$, а массовое число — $20$.

В результате захвата электрона и испускания протона образуется ядро ${}^{19}_{10}\mathrm{Ne},$ в котором содержится $10$ протонов и $9$ нейтронов.

Ответ: $109$.
}

% Source sdamgia-11276: https://phys-ege.sdamgia.ru/problem?id=11276
\BankTask{sdamgia-11276}{
Ядро ${}^{128}_{52}\mathrm{Te}$ может испытывать двойной бета-распад, при котором образуются два электрона, два антинейтрино и дочернее ядро (продукт распада). Сколько протонов и сколько нейтронов содержится в дочернем ядре?

\begin{center}\small\renewcommand{\arraystretch}{1.25}\begin{tabular}{|p{0.450\linewidth}|p{0.450\linewidth}|}\hline
Число протонов & Число нейтронов\\\hline
 & \\\hline
\end{tabular}\end{center}

В ответе запишите число протонов и нейтронов слитно без пробела.
}{
\smallskip\noindent\textbf{Решение}\par
В результате такой реакции заряд ядра увеличивается на $2$, а массовое число не изменяется. Таким образом, у дочернего ядра будет $54$ протона и $128 - 54 = 74$ нейтрона.

Ответ: $5474$.
}

% Source sdamgia-11674: https://phys-ege.sdamgia.ru/problem?id=11674
\BankTask{sdamgia-11674}{
Ядро изотопа углерода ${}^{14}_{6}\mathrm{C}$ претерпело электронный $\beta $-распад, в результате чего образовалось новое ядро X. Какой порядковый номер в таблице Д. И. Менделеева имеет соответствующий ядру X химический элемент и сколько нуклонов входит в состав ядра X?

\begin{center}\small\renewcommand{\arraystretch}{1.25}\begin{tabular}{|p{0.450\linewidth}|p{0.450\linewidth}|}\hline
Порядковый номер в таблице Д. И. Менделеева & Число нуклонов\\\hline
 & \\\hline
\end{tabular}\end{center}

В ответе запишите порядковый номер в таблице Д. И. Менделеева и число нуклонов слитно без пробела.
}{
\smallskip\noindent\textbf{Решение}\par
Запишем происходящую реакцию: ${}^{14}_{6}\mathrm{C} \rightarrow  {}^{A}_{Z}\mathrm{X} + _ - 1 ^{0} e.$ По закону сохранения массы: $A = 14$. По закону сохранения заряда: Z = $7$.

Ответ: $714$.
}

% Source sdamgia-25373: https://phys-ege.sdamgia.ru/problem?id=25373
\BankTask{sdamgia-25373}{
Ядро атома радия ${}^{228}_{88}\mathrm{Ra}$ претерпевает бета-распад. Какими станут массовое и зарядовое число образовавшегося ядра? Ответ запишите в виде двух чисел без пробела и запятой.
}{
\smallskip\noindent\textbf{Решение}\par
При бета-распаде вылетает электрон $^{0} _ - 1e.$ Запишем происходящую реакцию: ${}^{228}_{88}\mathrm{Ra} \rightarrow  {}^{228}_{89}\mathrm{Ac} + ^{0} _ - 1e + \overline\nu_e.$ Исходя из законов сохранения массы и заряда, массовое число останется равным $228$, зарядовое число увеличится на $1$ и станет равно $89$.

Ответ: $22889$.
}

% Source sdamgia-29992: https://phys-ege.sdamgia.ru/problem?id=29992
\BankTask{sdamgia-29992}{
Ядро ${}^{174}_{78}\mathrm{Pt}$ испытывает $\alpha $-распад, при этом образуются $\alpha $-частица и ядро химического элемента ${}^{A}_{Z}\mathrm{X} .$ Определите заряд Z (в единицах элементарного заряда) ядра X.
}{
\smallskip\noindent\textbf{Решение}\par
Альфа-частица обладает зарядом, равным $2$. По закону сохранения заряда при альфа-распаде $78 = 2$ + Z, откуда Z = $78 - 2 = 76$.

Ответ: $76$.
}

% Source sdamgia-38198: https://phys-ege.sdamgia.ru/problem?id=38198
\BankTask{sdamgia-38198}{
Ядро изотопа полония ${}^{195}_{84}\mathrm{Po}$ испытывает $\alpha $-распад, при этом образуются $\alpha $-частица и ядро химического элемента ${}^{A}_{Z}\mathrm{X}.$ Каков заряд Z ядра X (в единицах элементарного заряда)?
}{
\smallskip\noindent\textbf{Решение}\par
При альфа-распаде зарядовое число уменьшается на $2$. Следовательно, заряд ядра Х равен $82$.

Ответ: $82$.
}


% Source A2B69D: https://ege.fipi.ru/bank/questions.php?search=1&proj=BA1F39653304A5B041B656915DC36B38&qid=A2B69D
\BankTask{A2B69D}{
Радиоактивный изотоп висмута ${}_{83}^{212}\mathrm{Bi}$ претерпевает $\alpha$-распад. Укажите число протонов и число нейтронов в образовавшемся ядре.

\begin{center}
\begin{tabular}{|c|c|}
\hline
Число протонов & Число нейтронов \\
\hline
\rule{0pt}{1.5em} & \\ \hline
\end{tabular}
\end{center}
}{
\smallskip\noindent\textbf{Решение от GPT}\par
При $\alpha$-распаде массовое число уменьшается на $4$, а зарядовое — на $2$: $A'=212-4=208$, $Z'=83-2=81$. Поэтому в дочернем ядре $81$ протон и $N'=208-81=127$ нейтронов. Ответ: $81;\ 127$.
}

% Source 2B6FF4: https://ege.fipi.ru/bank/questions.php?search=1&proj=BA1F39653304A5B041B656915DC36B38&qid=2B6FF4
\BankTask{2B6FF4}{
Ядро таллия ${}_{81}^{210}\mathrm{Tl}$ испытывает электронный $\beta$-распад, при этом образуются электрон и ядро элемента ${}_Z^A X$. Каковы заряд $Z$ (в единицах элементарного заряда) и массовое число $A$ ядра $X$?

\begin{center}
\begin{tabular}{|c|c|}
\hline
Заряд ядра $Z$ & Массовое число ядра $A$ \\
\hline
\rule{0pt}{1.5em} & \\ \hline
\end{tabular}
\end{center}
}{
\smallskip\noindent\textbf{Решение от GPT}\par
При электронном $\beta^-$-распаде нейтрон превращается в протон с испусканием электрона и электронного антинейтрино. Поэтому зарядовое число увеличивается на единицу, а массовое сохраняется: $Z=81+1=82$, $A=210$. Ответ: $82;\ 210$.
}

% Source AB2CB2: https://ege.fipi.ru/bank/questions.php?search=1&proj=BA1F39653304A5B041B656915DC36B38&qid=AB2CB2
\BankTask{AB2CB2}{
Образец радиоактивного радия \({}_{88}^{224}\mathrm{Ra}\) находится в закрытом сосуде, из которого откачан воздух. Ядра радия испытывают \(\alpha\)-распад с периодом полураспада \(3{,}6\,\text{суток}\). Определите число моль гелия в сосуде через \(10{,}8\,\text{суток}\), если образец в момент его помещения в сосуд имел в своём составе \(6\cdot10^{23}\) атомов радия-224, а атомов гелия в сосуде не было.
}{
\smallskip\noindent\textbf{Решение от GPT}\par
За время \(t=10{,}8\,\text{суток}\) проходит \(t/T=10{,}8/3{,}6=3\) периода полураспада, поэтому распадается \(N_0(1-2^{-3})=\frac78N_0\) ядер радия. При каждом \(\alpha\)-распаде образуется одно ядро гелия, поэтому \(\nu_{\mathrm{He}}=\frac{7N_0}{8N_A}=\frac{7\cdot6\cdot10^{23}}{8\cdot6\cdot10^{23}}=0{,}875\ \text{моль}\). Ответ: \(0{,}875\,\text{моль}\).
}

% Source 38C6C8: https://ege.fipi.ru/bank/questions.php?search=1&proj=BA1F39653304A5B041B656915DC36B38&qid=38C6C8
\BankTask{38C6C8}{
Образец радиоактивного полония ${}_{84}^{218}\mathrm{Po}$ находится в закрытом сосуде, из которого откачан воздух. Ядра полония испытывают $\alpha$-распад с периодом полураспада $3\,\text{мин}$. Определите число моль полония-218 в сосуде через $9\,\text{мин.}$, если образец в момент его помещения в сосуд имел в своём составе $2{,}4\times10^{23}$ атомов полония-218.
}{
\smallskip\noindent\textbf{Решение от GPT}\par
За время $9\,\text{мин}$ проходит $t/T=9/3=3$ периода полураспада. Число оставшихся атомов полония-218 равно $N=N_0\cdot2^{-3}=3\times10^{22}$. Количество вещества $\nu=N/N_A=\dfrac{3\times10^{22}}{6\times10^{23}}\,\text{моль}=0{,}05\,\text{моль}$. Ответ: $0{,}05\,\text{моль}$.
}

% Source BD7F43: https://ege.fipi.ru/bank/questions.php?search=1&proj=BA1F39653304A5B041B656915DC36B38&qid=BD7F43
\BankTask{BD7F43}{
Установите соответствие между видами радиоактивного распада и уравнениями, описывающими этот процесс.

\TaskMatchingColumns[.25\linewidth]{Виды распада}{
А) альфа-распад\par
Б) бета-распад
}{Уравнения}{
1) \({}^{238}_{92}\mathrm{U}+{}^{22}_{10}\mathrm{Ne}\to{}^{256}_{102}\mathrm{No}+4{}^{1}_{0}\mathrm{n}\)\par
2) \({}^{238}_{93}\mathrm{Np}\to{}^{238}_{94}\mathrm{Pu}+{}^{0}_{-1}\mathrm{e}+\widetilde{\nu}_e\)\par
3) \({}^{227}_{89}\mathrm{Ac}\to{}^{223}_{87}\mathrm{Fr}+{}^{4}_{2}\mathrm{He}\)\par
4) \({}^{209}_{83}\mathrm{Bi}+{}^{2}_{1}\mathrm{H}\to{}^{105}_{43}\mathrm{Tc}+{}^{102}_{41}\mathrm{Nb}+4{}^{1}_{0}\mathrm{n}\)
}

\medskip
\begin{tabular}{|c|c|}
\hline
А & Б \\
\hline
\phantom{3} & \phantom{2} \\
\hline
\end{tabular}
}{
\smallskip\noindent\textbf{Решение от GPT}\par
При альфа-распаде испускается ядро \({}^{4}_{2}\mathrm{He}\): массовое число уменьшается на 4, зарядовое — на 2; это уравнение 3. При электронном бета-распаде испускаются электрон и электронное антинейтрино, массовое число сохраняется, а зарядовое увеличивается на 1; это уравнение 2. Уравнения 1 и 4 описывают вынужденные ядерные реакции с налетающей частицей. Ответ: А — 3, Б — 2.
}

% Source FF6C82: https://ege.fipi.ru/bank/questions.php?search=1&proj=BA1F39653304A5B041B656915DC36B38&qid=FF6C82
\BankTask{FF6C82}{
Выберите среди приведённых во втором столбце ядерных реакций те, которые являются примерами реакций альфа- и бета-распада.

\TaskMatchingColumns[.25\linewidth]{Вид ядерной реакции}{
А) альфа-распад\par
Б) бета-распад
}{Ядерные реакции}{
1) ${}_{77}^{176}\mathrm{Ir}\to{}_{75}^{172}\mathrm{Re}+{}_{2}^{4}\mathrm{He}$\par
2) ${}_{92}^{238}\mathrm{U}+{}_{10}^{22}\mathrm{Ne}\to{}_{102}^{256}\mathrm{No}+4{}_{0}^{1}n$\par
3) ${}_{48}^{113}\mathrm{Cd}+{}_{0}^{1}n\to{}_{48}^{114}\mathrm{Cd}+\gamma$\par
4) ${}_{71}^{178}\mathrm{Lu}\to{}_{72}^{178}\mathrm{Hf}+{}_{-1}^{0}e+\tilde{\nu}_e$
}

\begin{tabular}{|c|c|}
\hline
А & Б \\
\hline
\phantom{0} & \phantom{0} \\
\hline
\end{tabular}
}{
\smallskip\noindent\textbf{Решение от GPT}\par
Реакция 1 является альфа-распадом: испускается ${}_{2}^{4}\mathrm{He}$, поэтому $A$ уменьшается на 4, а $Z$ — на 2. Реакция 4 является электронным бета-распадом: испускаются электрон и антинейтрино, $A$ сохраняется, а $Z$ увеличивается на 1. Реакция 2 — слияние ядер с испусканием нейтронов, реакция 3 — захват нейтрона с испусканием гамма-кванта. Ответ: А — 1, Б — 4.
}

\end{document}
